\section{Finite packet repayment and closure}
\label{sec:packet-repayment}

We complete the proof of Theorem~\ref{thm:packet-growth} by a finite
backward iteration. A cut of length $2h$ records the loss from coherent
superposition as fine velocity information. At a growth rate
$\betaq>\gamma/2$, the repayment estimate forces this quantity to
increase at each cut by $(4\betaq-2\gamma)h$, up to the finite errors
specified below, while its capacity is at most $2h$ up to those errors.
A sufficiently long finite chain is therefore impossible.
The information carried between cuts is encoded by a uniform
conditional probability bound, so it can be retained through the
subsequent branch selections.

Write \(\ell=\log(1/\eps)\) and divide all entropies by \(\ell\).
We use the finite budget of Section~\ref{sec:packet-composition}:
a conditional list of size \(\eps^{-\alpha}\) costs \(\alpha\),
and a fixed factor \(C\) costs \((\log C)/\ell\).  Depth bins have
width \(\xi\), and packet tails are cut at distance \(\eps^{-\theta}\)
in earlier packet units, with \(\theta<h/10\) at fixed depths and
\(h>0\).  Constants \(C_*\) count finitely many geometric and
entropy losses, independently of the aspect exponents and cut
locations.  Fixed phase and amplitude constants enter through
their logarithms divided by \(\ell\).  We may increase \(C_*\)
a finite number of times.

Figure~\ref{fig:packet-cells} explains the mismatch between the earlier
packet's transported uncertainty and the output position precision.
The fine velocity information below accounts for this mismatch.
\begin{figure}[tbp]
\centering
\begin{tikzpicture}[font=\small,>=Stealth,
  dimension/.style={<->,line width=0.45pt}]
  \draw[->,figuregray!70] (-3.05,0) -- (4.18,0);
  \node[anchor=north east] at (4.18,-0.71) {velocity};
  \draw[->,figuregray!70] (0,-2.02) -- (0,2.50)
    node[above,text=black] {position at $C_t$};
  \filldraw[fill=figuregray!9,draw=figuregray,line width=0.85pt]
    (-2.2,-0.5) rectangle (2.2,0.5);
  \filldraw[fill=figureblue!17,draw=figureblue,line width=0.85pt]
    (-0.73333,-1.5) rectangle (0.73333,1.5);
  \draw[figuregray,dashed,line width=0.45pt]
    (-0.73333,-0.5) -- (-0.73333,0.5);
  \draw[figuregray,dashed,line width=0.45pt]
    (0.73333,-0.5) -- (0.73333,0.5);
  \draw[dimension] (-0.73333,1.78) -- (0.73333,1.78)
    node[midway,above=2pt,text=figureblue,fill=white,inner sep=1pt] {$\eps^{u_t+h}$};
  \draw[dimension] (-2.2,-2.13) -- (2.2,-2.13)
    node[midway,below=2pt] {$\eps^{u_t}$};
  \draw[dimension] (-2.63,-1.5) -- (-2.63,1.5)
    node[midway,left=5pt,text=figureblue,fill=white,inner sep=1pt] {$\eps^{u_t+t-h}$};
  \draw[dimension] (2.61,-0.5) -- (2.61,0.5)
    node[midway,right=5pt,fill=white,inner sep=1pt] {$\eps^{u_t+t}$};
  \draw[figureblue,line width=0.85pt] (-3.4,-2.98) -- (-2.92,-2.98);
  \node[anchor=west,font=\footnotesize] at (-2.81,-2.98)
    {earlier packet, transported from $r=t-2h$};
  \draw[figuregray,line width=0.85pt] (-3.4,-3.35) -- (-2.92,-3.35);
  \node[anchor=west,font=\footnotesize] at (-2.81,-3.35)
    {output packet at $t$};
\end{tikzpicture}
\caption{One coordinate pair of packet uncertainty cells at a common
output anchor, with $0<h\le t/2$, $t\le\Planck$ and
$u_t=(\Planck-t)/2$.
The earlier packet has $u_r=u_t+h$: its velocity cell is finer by depth
$h$, while its transported position uncertainty is coarser by depth $h$.
Both displayed rectangles have the same area,
$\eps^{u_t+h}\eps^{u_t+t-h}
=\eps^{u_t}\eps^{u_t+t}=\eps^{\Planck}$.
The rectangles are drawn concentrically to compare their scales and
represent comparable widths after transport; fixed factors, shear and
subpower tail enlargements are suppressed.
An earlier packet can therefore contribute across several fine output
position cells.  Its transported center alone does not determine the
fine output position label.}
\label{fig:packet-cells}
\end{figure}


\subsection{The incoming fine-velocity information}

Here is the precise meaning of the velocity information carried between
steps.  At time \(t\), let
\[
 u_t=(\Planck-t)/2,\qquad
 P_f=(C_t,\text{isotropic velocity and position cells of width }
                  \eps^{u_t-h}).
\]
The position width in this formula is \(\eps^{t+u_t-h}\).
Fixed dilations and translations of the grids are allowed.  Additional
tags are allowed only if their conditional number, given the geometric
part of \(P_f\), is at most \(\eps^{-\alpha}\); their cost is included
in \(\alpha\).

In the elongated case \(n_f(P_f)\) is a unit normal, modulo sign, and
the normal of the aperture \(D\) is within \(C\eps^h\) of \(n_f(P_f)\).
The label \(W_f\) locates the velocity component along \(n_f(P_f)\)
in intervals of length \(\eps^{u_t}\).  In the isotropic case \(W_f\)
locates both velocity components at that precision.  An incoming
credit \(k_f\geq0\), with error \(\alpha_f\), means that on the
normalized law being used,
\begin{equation}
 \sup_w\PP(W_f=w\mid P_f=p)\leq
       \eps^{\,k_f-\alpha_f}
       \quad\text{for every \(p\) of positive probability}.
 \label{eq:incoming-packet-credit}
\end{equation}
In particular \(\Ent(W_f\mid P_f)\geq k_f-\alpha_f\).
On the first step we set \(k_f=0\) and impose no incoming direction
or fine-velocity condition.

We will use the following exact retention rule.  It specifies how
\eqref{eq:incoming-packet-credit} is maintained when selecting branches.

\begin{lemma}[Retention of a conditional probability bound]
\label{lem:packet-credit-retention}
Suppose a finite law satisfies
\(\PP(W=w\mid P=p)\leq\eps^{k-\alpha}\).
Let \(A\) have probability \(p_A>0\), and let \(\zeta>0\).
Under the law conditioned on \(A\), the set of fibers on which
\[
 \PP(A\mid P)<p_A\eps^\zeta
\]
has probability at most \(\eps^\zeta\).  On every remaining fiber,
\begin{equation}
 \PP(W=w\mid P,A)\leq
   \eps^{\,k-\alpha-\zeta-\logeps{1/p_A}}.
 \label{eq:credit-retention}
\end{equation}
Removing the exceptional fibers does not change this conditional
bound.  The same assertion holds when \(P\) includes other previously
fixed labels.
\end{lemma}

\begin{proof}
The exceptional probability, in the conditioned law, is
\[
 p_A^{-1}\sum_{\{p:\,\PP(A\mid P=p)<p_A\eps^\zeta\}}
       \PP(P=p)\PP(A\mid P=p)\leq\eps^\zeta.
\]
On a remaining fiber divide
\(\PP(W=w,A\mid P=p)\leq\PP(W=w\mid P=p)\)
by \(\PP(A\mid P=p)\).  Conditioning subsequently on a set of \(P\)
values does not alter the conditional law given \(P\).
\end{proof}

Every removal below is also a score split under
Section~\ref{sec:packet-composition}; its probability alone does
not control the remaining score.


\subsection{Constructing the finite cut}

A prepared cut starts with a boundary step from \(r=t-2h>0\) to \(t\),
with \(u_t>h\),
as supplied by Proposition~\ref{prop:boundary-alternatives}.
The current endpoint law will be denoted by \(\pi^{\rm in}\).
It is a probability law on a finite sample space \(\Omega\).
A sample \(\omega\) records \(O(\omega),D(\omega)\), all incoming
labels \(P_f,W_f\) and their extra tags when present, and any auxiliary
randomness already used to choose those labels.  In particular, an
incoming label need not be a deterministic function of \((O,D)\).
For any endpoint law \(\nu\) on \(\Omega\), write \(P_t(\nu)\)
for the score of the original input with its original masses.
If \(k_f>0\), the bound \eqref{eq:incoming-packet-credit} holds under
\(\pi^{\rm in}\), with its current error.  On the first step
\(k_f=0\).

The finite geometry is the following.  The endpoint shape has
\(a=u_t\) to depth error at most \(\eta\).  In the isotropic
alternative the apertures are isotropic and \(b_*=0\).  In the
elongated alternative \(b_*=b\), the endpoint aspect \(e=a-b\)
is at least \(e_0>0\).  In that case require
\begin{equation}
 C_*(h+\eta)<\gamma e_0 .
 \label{eq:repayment-eccentricity}
\end{equation}
The grids and aperture menu admit all coarsenings below.  We take
\(C_*\ge2\).  Thus in the elongated case
\(b\le u_t+\eta-e_0\le u_t-h\), after increasing \(C_*\) if
necessary; in the isotropic case \(b_*=0<u_t-h\).  We will use
\begin{equation}
 b_*\le u_t-h
 \label{eq:repayment-localizer-depth}
\end{equation}
when comparing the incoming coarse label with the localizer.

Use the fixed analytic, domain, depth, grid, and count bounds of
Section~\ref{sec:packet-composition}.  For supremum masses, fix at
each site a test within a fixed factor of the supremum and call its
mass \(m_O(f)\).  This changes any endpoint score by at most
\(C/\ell\), included in \(\mathcal R\), and the Bessel bounds are
uniform in these choices.  Apply the floor
\(m_O(f)\ge\eps^{A_0}\|f\|_2^2\) of
\eqref{eq:packet-output-floor} and the fixed-shape score splits,
recording their probabilities for credit retention.

Fix the depths, then \(\theta<h/10\) and \(\xi\).  By
\eqref{eq:repayment-localizer-depth},
\(\theta<(u_r-b_*)/4\), as required by
Lemma~\ref{lem:packet-decomposition}.  Choose its amplitude error
\(C_A\eps^A\|f\|_2\) with \(A>A_0/2+1\).  If \(Q_0\) bounds
all relevant count exponents, choose the coefficient cutoff
\begin{equation}
 B>2A+3Q_0+10 .
 \label{eq:repayment-coefficient-cutoff}
\end{equation}
uniformly over geometric localizers.  The later transport order
may be increased after this coefficient family is fixed.

For one geometric localizer \(J_0\), let \(\iota\) range over its
coefficient bins and greedy cap classes.  The associated whole class
is a finite set \(\mathcal T_\iota^{\rm all}\) of earlier atoms,
with coefficient mass bin \(m(\iota)\), cap maximum bin \(g(\iota)\),
and assigned cap \(E_\iota(T)\).  Lemma~\ref{lem:cap-class} gives
the cap upper bound on this set and the lower bound on each of its
whole assigned fibers.  At this point \(\iota\) indexes available
classes; no class has yet been assigned to an endpoint sample as its
final winner.  The buffer used next is chosen for all possible
contributors in this finite family, before any of its classes is
subdivided.

During preparation, a replacement loss is the increase in score
when the aperture is changed at fixed masses, evaluated on the
specified current marginal.  Section~\ref{sec:packet-composition}
supplies these losses from the current deficits and charges each
split.  This also applies to the two-copy marginals in the next
lemma, since adjoining the unsplit conditional copies preserves
the old marginal.

For every possible contributor, the velocity difference from its
output center is \(O(\eps^{u_t})\).  After referring the positions
to \(C_r\), their difference is
\[
 O(\eps^{r+u_t})+
 O(\eps^{-\theta}\eps^{r+u_r})
       =O(\eps^{r+u_t}),
 \qquad u_r=u_t+h,\quad \theta<h/10.
\]
Use isotropic cells with side \(L_{\rm grid}\) times
\(\eps^{u_t}\) in velocity and \(\eps^{r+u_t}\) in position,
where \(L_{\rm grid}\) is a sufficiently large fixed constant.
A uniformly random translation places any output center farther
than these possible differences from every face with probability
at least \(3/4\).  Averaging over translations therefore gives
one translation for which a part of at least that probability has
the coarse label
\[
 P=(C_r,V_{u_t},(X+(C_r-C_t)V)_{r+u_t})
\]
in common with every possible contributing earlier atom.
Retain this part as a score split and denote its normalized
endpoint law by \(\pi^{\rm buf}\).  Its probability and deficit
are charged by the finite selection rule above.  This completes
the geometric preparation in the isotropic case.

\begin{lemma}[Prediction of a normal from isotropic replacements]
\label{lem:orientation-prediction}
Assume the elongated geometry.  On \(\pi^{\rm buf}\), suppose the
replacement losses for the isotropic apertures of widths \(a\) and
\(b\) are included in \(C_*\eta\).  Take two copies
\((O,D),(O',D')\), independent conditional on \(P\), and set
\[
 A_{\rm sep}=\{\angle(n(D),n(D'))\ge\eps^{4h}\}.
\]
If this event has probability at least \(1/4\), require the same
replacement bounds on both marginals of the law conditioned on
\(A_{\rm sep}\).  Then there is a
deterministic normal \(n(P)\), defined on every coarse cell used
by the atom partition, such that
\[
 \angle(n(D),n(P))\le C\eps^{4h}
\]
on a part of \(\pi^{\rm buf}\)-probability at least \(3/4\).
\end{lemma}

\begin{proof}
Let \(e=a-b\), and let \(B_P\) be the isotropic coarsening of \(P\)
to width \(b\).  Grid list errors will be included in \(C_*\eta\).
Only \(2h\) extra depth for time and \(2h\) for each position
coordinate are missing from \(P\) in order to specify \(O\), so
\begin{equation}
 \Ent(O\mid P)\leq6h+C_*\eta.
 \label{eq:prediction-missing-data}
\end{equation}
Replacing \(D\) by the isotropic aperture of width \(a\) gives
\begin{equation}
 \Ent(P\mid D)\geq(1+\gamma)e-6h-C_*\eta.
 \label{eq:prediction-lower}
\end{equation}
Indeed the weight increases by \((1+\gamma)e\), and the new
conditional observation entropy is nonnegative.  Replacing \(D\)
by the isotropic width-\(b\) aperture at \(t\), which has at most
\(6h+C_*\eta\) information beyond \(B_P\), gives
\begin{equation}
 \Mut(P;D\mid B_P)\leq(1-\gamma)e+6h+C_*\eta.
 \label{eq:prediction-upper}
\end{equation}
Both inequalities are comparisons on the marginal law on which
they are invoked.

Suppose \(A_{\rm sep}\) has probability at least \(1/4\), and
condition the two-copy law on this event.
For the resulting law, conditional dependence of the two copies
given \(P\) is at most
\(\log(1/\PP(A_{\rm sep}))/\ell\leq(\log4)/\ell\).
This follows by comparing its density with the original conditional
product and using the chain rule for relative entropy.
The marginal replacement estimates
\eqref{eq:prediction-lower}--\eqref{eq:prediction-upper}
continue to hold with their charged conditional deficits.

Each aperture, referred to \(C_r\), constrains both columns to a
normal strip of width a fixed multiple of \(\eps^{a-C_*\eta}\).
Two such strips whose normals have angle at least \(\eps^{4h}\)
have intersection of diameter at most
\(C\eps^{a-4h-C_*\eta}\).
Covering both columns by \(P\) cells consequently gives
\begin{equation}
 \Ent(P\mid D,D')\leq16h+C_*\eta.
 \label{eq:prediction-intersection}
\end{equation}
Here a strip intersection is used separately for velocity and root
position, each in two dimensions; hence \(4(4h)=16h\).
The coarse label \(B_P\) has bounded ambiguity given either aperture.
The identity
\[
 \Mut(P;D'\mid D,B_P)-\Mut(P;D'\mid B_P)
 =\Mut(D;D'\mid P,B_P)-\Mut(D;D'\mid B_P)
\]
and conditional product comparison therefore imply that its left
side is at most \(C_*\eta\).
The first term is at least
\((1+\gamma)e-22h-C_*\eta\), by
\eqref{eq:prediction-lower} and
\eqref{eq:prediction-intersection}; the second is at most
\((1-\gamma)e+6h+C_*\eta\), by
\eqref{eq:prediction-upper} for the other marginal.
Thus \(2\gamma e\leq28h+C_*\eta\), contrary to the eccentricity
hypothesis after fixing \(C_*\).
It follows that the complementary angle-proximity event has
probability at least \(3/4\).  For each \(P=p\), choose one normal
maximizing the conditional mass of a ball of radius
\(C\eps^{4h}\).  Averaging the two-copy probability shows that the
single-copy part predicted by these normals has mass at least
\(3/4\).

On coarse cells of zero probability choose an arbitrary normal.
This extends the table \(n(P)\) to every cell that will index an atom
subdivision, without changing the asserted event.
\end{proof}


\subsection{The subclasses and the current endpoint law}

Retain the buffered and, in the elongated case, predicted-normal
parts just constructed.  The table \(n(P)\) and the translated grid
are now fixed.  Their selection probabilities, together with the
earlier mass-floor and shape selections, are included in one
preparation event \(A_{\rm pre}\) in the original sample space
\(\Omega\).  We will restore the incoming conditional probability
bound after the atom partition and the winning map have been fixed.

For an earlier atom \(T\), let \(P(T)\) be its buffered coarse cell.
In the elongated case first mark the atom according to whether the
normal of its assigned aperture \(E_\iota(T)\) lies within
\(2C\eps^h\) of \(n(P(T))\).  This divides each whole class into
aligned and misaligned atoms; both parts remain available when the
winner is chosen.  In the isotropic case no such mark is needed.

Let \(W(T)\) record the velocity at depth \(u_r=u_t+h\).  It records
the component along \(n(P(T))\) in the elongated case and both
components in the isotropic case.  These are the original depth
units; after subtracting \(b_*\), the fine spacing is
\(\eps^{u_r-b_*}\).  Within each marked part of
\(\mathcal T_\iota^{\rm all}\), and for each \(P\), bin
\[
 \logeps{\#\{T:(P(T),W(T))=(p,w)\}}
\]
with mesh \(\xi\).  Among fibers in one such size bin, bin the
logarithm of the number of represented \(W\)'s at \(P\), again with
mesh \(\xi\).  Write \(k\ge0\) for this second numerical tag and
\(d=1\) in the elongated case, \(d=2\) in the isotropic case.
For each resulting subclass and every represented \(P\),
\begin{equation}
 \eps^{-k+\xi}\leq \#\{W\mid P\}\leq\eps^{-k-\xi},
 \qquad k\leq dh+C\xi+\frac{\log C}{\ell}.
 \label{eq:velocity-fiber-regularity}
\end{equation}
The last bound covers a velocity cell of width \(C\eps^{u_t}\)
by cells of width \(\eps^{u_r}\) in \(d\) dimensions.
The sizes of two nonempty \((P,W)\) fibers in this subclass differ
by at most \(\eps^{-\xi}\).  Hence its uniform atom law satisfies
\begin{equation}
 \PP(W=w\mid P,J)\leq\eps^{k-2\xi}.
 \label{eq:prepared-velocity-probability}
\end{equation}
Here \(J\) will be the label of the subclass when it is chosen.
The subdivision groups the numerical fiber-size and \(k\) tags
across different \(P\)'s; \(P\) itself remains random in the
subclass and is not appended as an auxiliary value of \(J\).
This distinction keeps the later conditional cost of \(J\) small.

These are the final subdivisions that change atom sets.  Later tags,
such as total size, label a whole subclass.  Let \(S_\eps\) bound
the number of available subclasses at a geometric localizer; at
fixed depths \(\log S_\eps=o(\ell)\).  Their labels contain \(J_0\),
the \(\iota\) tags, the alignment mark when present, the two
numerical velocity tags, and any whole-subclass tags.  The winner map will choose one
such tuple \(J(\omega)\).  For its atom set \(\mathcal T_J\), set
\(\mathcal T_J^{\rm all}:=\mathcal T_\iota^{\rm all}\),
\(E(T):=E_\iota(T)\), and
\[
 N(J)=\#\mathcal T_J,\qquad
 N_{\rm all}(J)=\#\mathcal T_\iota^{\rm all},\qquad
 m(J)=m(\iota),\qquad g(J)=g(\iota).
\]
Thus \(N_{\rm all}\) still counts the whole greedy class before
alignment and velocity subdivision.  Its cap upper bound persists
on \(\mathcal T_J\); its assigned-fiber lower bound continues to
refer to the whole class.

Write \(\mathfrak a_O(\varphi)\) for the normalized pairing with
the fixed site test, so \(m_O(\varphi)=|\mathfrak a_O(\varphi)|^2\),
and \(h_J=\sum_{T\in\mathcal T_J}c_Tg_T\).  Put
\[
 q_0=C_A\eps^{A-A_0/2}<\tfrac12.
\]
Choose at each retained sample a subclass maximizing
\(|\mathfrak a_O(h_J)|\), breaking ties by a fixed rule.  All atom
sets, superpositions, and this map to \(J\) are now fixed.  Applied
to these final subclasses, \eqref{eq:packet-winner-mass} gives
\begin{align}
 |\mathfrak a_O(h_J)|
 &\ge\frac{1-q_0}{S_\eps}\sqrt{m_O(f)}
 \ge\frac{1-q_0}{S_\eps}\eps^{A_0/2}\|f\|_2,
 \label{eq:repayment-winner-floor}\\
 m_O(f)&\le (1-q_0)^{-2}S_\eps^2m_O(h_J)
          \le4S_\eps^2m_O(h_J).
 \label{eq:repayment-original-winning}
\end{align}
The mass-term loss
\(3\logeps{S_\eps}+3\logeps{1/(1-q_0)}\) is included in
\(\mathcal R\).  In the elongated case both alignment marks
participated in this choice.

After this finite menu, winner map, and coefficient cutoff are fixed,
Lemma~\ref{lem:packet-decomposition} allows one common transport power
\begin{equation}
 A_{\rm tr}>(B+Q_0+1)/2
 \label{eq:repayment-transport-order}
\end{equation}
for every retained subclass, without a new coefficient cutoff.
Its constants depend only on the fixed data and \(Q_0\).  This
order controls the alignment collapse and the later candidate
truncation, and is unchanged by subsequent subclass selections.

Take the endpoint law conditioned on \(A_{\rm pre}\).  If \(k_f>0\),
apply Lemma~\ref{lem:packet-credit-retention} to this event, using
its probability under \(\pi^{\rm in}\), and remove the exceptional
\(P_f\) fibers.
This changes the endpoint law and its full \(J\) marginal, while
leaving the atom sets, counts, \(g,m\), and winner map fixed.
The pointwise geometric and mass comparisons persist.  Apply
Proposition~\ref{prop:finite-packet-composition} anew to this
law; its new deficit inequality controls the mean
\(\EE\logeps{N_{\rm all}/N}\) under its new \(J\) marginal.

For each later numerical-tag or alignment restriction \(A\), use
its probability \(p_A\) in the current normalized law.  When
\(k_f>0\), repeat the retention lemma and the exceptional-fiber
score split.  In either case reapply composition with the same
frozen data.  Each retention application adds
\begin{equation}
 \zeta+\logeps{1/p_A}
 \label{eq:repayment-credit-cost}
\end{equation}
to the incoming error.  Use the weighted tag rule of
Section~\ref{sec:packet-composition} with its current union mass
\(p\) and shifted deficit \(R\): it supplies probability at least
\(p/(4S_\eps)\) and
conditional deficit at most \(4R/p\).  Only a prescribed finite
number of these operations is used.

The following exclusion determines which marked subclasses can remain.

\begin{lemma}[Alignment of the selected earlier aperture]
\label{lem:orientation-alignment}
For the elongated prepared cut, let \(\nu\) be a current endpoint
law supported on \(A_{\rm pre}\),
with the frozen winner map, and let \(d_i,R_t\) be its deficits and
bound from Proposition~\ref{prop:finite-packet-composition}.
Let \(B\) be a \(J\)-measurable event of misaligned winners, with
\(p=\nu(B)>0\), and put \(\varrho=\nu(\,\cdot\mid B)\).
In local coordinates write \(e=a-b\) and \((j,l)\) for the
assigned-cap widths.  Assume the common shape bins satisfy
\[
 l=e+h+O(\eta),\qquad l-j\ge e_0-O(\eta),
\]
as supplied by the boundary alternative, and that the prepared
packet, cap, and label errors are included in \(C_*\eta\).
Then
\begin{equation}
 p\,[2\gamma e_0-C_*(h+\eta)]
 \le \EE_\nu\!\left[\boldsymbol1_B(d_1+3d_2+3d_3)\right]
 \le 3R_t .
 \label{eq:alignment-misaligned-mass}
\end{equation}
Thus a positive bracket bounds the current probability of
misaligned winners by the current finite deficit.
\end{lemma}

\begin{proof}
Because \(B\) is \(J\)-measurable,
\(\varrho(\,\cdot\mid J)=\nu(\,\cdot\mid J)\) on \(B\), so the
local deficits are inherited.  Work at a fixed \(J\), subtract
\(b_*=b\) from widths, put \(\delta_J=d_2(J)+d_3(J)\), and let
\(M_j\) be the earlier isotropic depth-\(j\) cell.

For this comparison only, let
\(\mathcal O_J^\varrho=\{O(\omega):\varrho(\omega)>0,\ J(\omega)=J\}\).
For fixed \(P,C_t\), collapse the winning masses as
\[
 m^{\rm buf}_{P,C_t}
 =\sum_{\substack{O\in\mathcal O_J^\varrho\\O\mapsto(P,C_t)}}m_O(h_J).
\]
This sum counts each geometric site once, with the allowed bounded
lattice multiplicity.  Every tested site is buffered, so each of
its possible contributors shares its \(P\).  Write
\(\mathcal T_{J,P}=\{T\in\mathcal T_J:P(T)=P\}\) and
\(N_{J,P}=\#\mathcal T_{J,P}\).  The finite mass bound of
Proposition~\ref{prop:packet-trivial-growth}, with the fixed floor
and tail comparison, bounds the collapsed mass by \(mN_{J,P}\)
with the stated finite error.  Only then sample a contributor:
\[
 \varrho^P(\omega,T)=\varrho(\omega)
   \frac{\boldsymbol1_{\{T\in\mathcal T_{J(\omega),P(\omega)}\}}}
        {N_{J(\omega),P(\omega)}}.
\]
The floor and tail comparison ensure \(N_{J,P}>0\) on the tested
samples.  This kernel preserves the whole endpoint sample; its
atom marginal may be biased.  The following entropies are
conditional on \(J\) under this law.

The entropy calculation of
Proposition~\ref{prop:packet-trivial-growth} bounds \(P_J\) above by
\[
 \Ent(T)+\tfrac32\logeps{m}+
 \tfrac12\{\Ent(T\mid D)+(1-\gamma)e\}
 +\Ent(C_t\mid T)+C_*\eta.
\]
Here \(\Ent(T)\le\logeps{N}\) and
\(\Ent(C_t\mid T)\le2h+C_*\eta\).  For every sampled atom at \(P\)
and every aperture \(D\) over that \(P\), the buffered velocity and
root-position columns differ from those of the paired output by
at most the cell widths.  Since \(a=u_t+O(\eta)\), a fixed
enlargement of the root-referred cap of \(D\) contains all these
atoms, with the charged depth error.  The cap upper bound gives
\(\Ent(T\mid D)\le g-(1-\gamma)e+C_*\eta\).
The definitions of \(d_2,d_3\) give
\[
 P_J=A_J+2\betaq h+2\mathcal R-\delta_J\ge A_J-\delta_J.
\]
Comparing the two score bounds yields
\begin{equation}
 \Ent(T)\ge\logeps{N}-\delta_J-C_*(h+\eta),\qquad
 \Ent(T\mid D)\ge g-(1-\gamma)e-2\delta_J-C_*(h+\eta).
 \label{eq:alignment-biased-entropy}
\end{equation}

Before subdivision the number of distinct assigned \(E\) labels
is at most
\(\eps^{-\logeps{N_{\rm all}}+g-w(j,l)-C_*\eta}\), by the
whole-class fiber lower bound of Lemma~\ref{lem:cap-class}.
Since \(E\) is assigned deterministically to \(T\),
\begin{equation}
 \Ent(T\mid E)\ge
 g-w(j,l)-d_1(J)-\delta_J-C_*(h+\eta).
 \label{eq:alignment-assigned-entropy}
\end{equation}
The whole-fiber count has thus cost \(d_1(J)\) for the retained,
possibly biased law.

On \(A_{\rm pre}\), the normal of \(D\) is within \(C\eps^{4h}\)
of \(n(P)\).  A misaligned \(E\) therefore makes angle at least
\(C\eps^h\) with it.  Their two column intersections fit isotropic
caps of depth \(e-C_*h-C_*\eta\), so
\[
 \Ent(T\mid D,E)\le g-2e+C_*(h+\eta).
\]
Subtracting from \eqref{eq:alignment-assigned-entropy} and using
\(l=e+h+O(\eta)\) gives, for fixed \(J\),
\[
 \Mut(T;D\mid E)\ge
 (1+\gamma)(e-j)-d_1(J)-\delta_J-C_*(h+\eta).
\]
Average under \(\varrho_J\), absorbing the common-bin errors:
\begin{equation}
 \Mut_{\varrho^P}(T;D\mid E,J)\ge
 (1+\gamma)(e-j)-\EE_\varrho(d_1+\delta_J)-C_*(h+\eta).
 \label{eq:alignment-information-lower}
\end{equation}
As \(E\) is a function of \(T\) given \(J\), and determines \(M_j\) up to
the charged lists,
\[
 \Mut_{\varrho^P}(T;D\mid E,J)
 \le \Mut_{\varrho^P}(T;D\mid M_j,J)+C_*\eta.
\]

For the reverse estimate, coarsen \(D\) to widths \(0,j\), using
the center of \(M_j\).  Its normal projection is confined by the
coarsened strip, and the tangent width is zero-depth, so the
resulting labels form a charged list determined by \(D\).
Enlarge each rectangle by a fixed factor to contain the entire
\(M_j\) cell; its projections are bounded by its diagonal.
The enlargement depends only on the rectangle label.  Resample it
conditional on \((M_j,J)\), independently of \(T\), and call it
\(U\).  It is faithful for every atom in that cell, has weight
\((1-\gamma)j\), and the cap bound gives on each \(J\) fiber
\[
 \Mut(T;U)\ge \Ent(T)-g+(1-\gamma)j-C_*\eta.
\]
Adjoin these kernels under \(\varrho_J\), retaining the notation
\(\varrho^P\).  The joint law of the enlarged coarsening and \(M_j\) is
preserved on each fiber.  Since \(M_j\) is a function of \(T\) given \(J\),
\[
 \Mut_{\varrho^P}(M_j;D\mid J)
 \ge \Mut_{\varrho^P}(M_j;U\mid J)-C_*\eta
 =\Mut_{\varrho^P}(T;U\mid J)-C_*\eta.
\]
Average the cap bound and the second inequality in
\eqref{eq:alignment-biased-entropy}.  Combining them with this
display gives
\begin{equation}
 \Mut_{\varrho^P}(T;D\mid M_j,J)\le
 (1-\gamma)(e-j)+2\EE_\varrho\delta_J+C_*(h+\eta).
 \label{eq:alignment-information-upper}
\end{equation}
The lower and upper bounds imply
\[
 2\gamma(e-j)\le
 \EE_\varrho[d_1+3(d_2+d_3)]+C_*(h+\eta).
\]
The common shape bins give \(e-j\ge e_0-h-O(\eta)\).
Multiply by \(p\), use
\(p\EE_\varrho[\cdot]=\EE_\nu[\boldsymbol1_B\,\cdot]\), and use nonnegativity
of the \(d_i\) and \eqref{eq:finite-packet-deficits}.  This proves
\eqref{eq:alignment-misaligned-mass}.
\end{proof}

For the elongated application, first select a branch with qualifying
common earlier shape bins by the weighted rule, perform credit
retention when needed, and recompose.  Let \(\nu\) be this current
law and let \(B\) be all its misaligned winners.  This event is
\(J\)-measurable.  If \(\nu(B)=0\), the law is already aligned;
otherwise \eqref{eq:alignment-misaligned-mass} applies to its
entire current mass.  At sufficiently small current deficit,
the aligned complement has positive probability, tending to one
with the positive gap fixed.  Restrict to that complement, then
apply the current-law retention and recomposition rule.
In either alternative, apply the same rule to any remaining common
numerical-bin selections and recompose after the last restriction.  Denote
the resulting endpoint law on \(\Omega\) by \(\pi\), and now set
\(P_t=P_t(\pi)\), the score of the original input with its original masses.
Every final elongated winning class is aligned.
The buffered property, and prediction when present, persist pointwise.
The incoming bound is
\begin{equation}
 \sup_w\pi(W_f=w\mid P_f=p)
 \le \eps^{k_f-\alpha_f}
 \quad\text{for every \(p\) with \(\pi(P_f=p)>0\)},
 \label{eq:repayment-selected-credit}
\end{equation}
where \(\alpha_f\) includes the initial error and the sum of
\eqref{eq:repayment-credit-cost}.  When \(k_f=0\), this estimate
will be replaced by nonnegativity of conditional mutual information.

Every expectation in the rest of this cut uses the selected marginal
\(\pi_J\), unless a different law is explicitly displayed.  In
particular, set
\begin{align}
 d_1(J)&=\logeps{N_{\rm all}(J)/N(J)},\nonumber\\
 A_J&=\logeps{N(J)}+\tfrac32\logeps{m(J)}+\tfrac12g(J),\nonumber\\
 K_r^*(\pi)&=\Ent_\pi(J)+b_*+
   \EE_\pi\left[\logeps{N_{\rm all}(J)}
                  +\tfrac32\logeps{m(J)}+\tfrac12g(J)\right],
 \nonumber\\
 B_r&=\Ent_\pi(J)+b_*+\EE_\pi A_J
       =K_r^*(\pi)-\EE_\pi d_1 .
 \label{eq:repayment-benchmark}
\end{align}
The last equality is exact.  The number \(B_r\) is a benchmark
formed with retained counts; it is not asserted to be an earlier
score.
For the final composition law, put
\(\Delta_t=K_0+\betaq t-P_t\).  Its current \(d_0\) gives
\begin{equation}
 B_r+2\betaq h-P_t
 =\Delta_t+2\mathcal R-d_0-\EE_\pi d_1
 \le\Delta_t+2\mathcal R .
 \label{eq:repayment-input-bridge}
\end{equation}

The earlier whole-class law and the returned retained-class law are,
respectively,
\begin{equation}
 \mu^{\rm all}(J,T)=\pi_J(J)
   \frac{\boldsymbol1_{\{T\in\mathcal T_J^{\rm all}\}}}{N_{\rm all}(J)},
 \qquad
 \mu(J,T)=\pi_J(J)
   \frac{\boldsymbol1_{\{T\in\mathcal T_J\}}}{N(J)}.
 \label{eq:repayment-earlier-laws}
\end{equation}
Both use the selected marginal \(\pi_J\).  Their observations are
the sites of \(T\) at time \(r\), and their masses are coefficient
tests of the original input \(f\).  With assigned caps \(E(T)\),
finite composition gives
\begin{align}
 P_r^{\rm all}&\ge K_r^*(\pi)-\mathcal R,\nonumber\\
 P_r^{\rm ret}&\ge K_r^*(\pi)-\tfrac32\EE_\pi d_1-\mathcal R
       =B_r-\tfrac12\EE_\pi d_1-\mathcal R .
 \label{eq:repayment-earlier-scores}
\end{align}
In the elongated case \(\mu\) continues the aligned assigned cap.
In the isotropic case the same atom law uses the exact isotropic boundary aperture
of width \(u_r\).  The horizon and packet menu include this cap by
Proposition~\ref{prop:boundary-alternatives}.  With \(b_*=0\),
the isotropic score satisfies the separate bound
\begin{equation}
 P_r^{\rm iso}\ge
 B_r-\tfrac12\EE_\pi(g-2u_r)-C\mathcal R .
 \label{eq:repayment-isotropic-score}
\end{equation}
This uses weight \(2u_r\) and nonnegative conditional entropy,
independently of the assigned-fiber gain.  A charged selection fixed
the common \(g\) bin, and the final credit trim, when needed,
preserved its pointwise support before \(\pi\) was named.  Thus \(g-2u_r\) has
the required small error under the displayed expectation.

Finally, \eqref{eq:prepared-velocity-probability} applies under
\(\mu(\,\cdot\mid J)\).  Since the numerical \(k\) tags are common
bins under \(\pi_J\), averaging in \(J\) gives
\begin{equation}
 \mu(W=w\mid P=p)\le\eps^{k-C_*\eta}
 \quad\text{for every \(p\) with \(\mu(P=p)>0\)}.
 \label{eq:repayment-outgoing-credit}
\end{equation}
This statement holds for either earlier aperture assignment under
\(\mu\), which is uniform within each retained class.


\subsection{From winning masses to candidate masses}

Before stating the repayment estimate, we finish the finite comparison
of masses that its proof will use.  Work in the translated coordinates
of a selected localizer \(J\), put \(u=u_t-b_*\), and write
\[
 \rho=\eps^u,\qquad s=\eps^{u+h},\qquad
 \sigma=\eps^{u+2h}=s^2/\rho,\qquad \lambda_{\rm loc}=s^{-2}.
\]
The buffered \(P\) cell has width \(\rho\).  Let \(Q=Q(J,O)\) record
\(P,C_t\), and the position cell of side \(s\) at \(C_t\).
For each such \(Q\), let \(\mathcal T_{J,Q}\subset\mathcal T_J\)
consist of the atoms at this \(P\) whose transported positions are
within \(C_0 R s\) of that cell, where \(R=\eps^{-\theta}\)
and \(C_0\) is a fixed constant large enough for the transport and
cell-boundary enlargements.  Absorb \(C_0\) into the fixed constants
in the coherence estimates.
Write \(N_{J,Q}=\#\mathcal T_{J,Q}\) and
\[
 h_{J,Q}=\sum_{T\in\mathcal T_{J,Q}}c_Tg_T.
\]
Let \(Z=Z(J,O)\) add to \(Q\) the output position at precision
\(\sigma\) along \(n(P)\) in the elongated case and in both
directions in the isotropic case.  These labels are functions of the
geometric endpoint observation and the fixed grid and normal table;
they do not include the arbitrary extra labels carried by \(\omega\).
On the full sample space write \(Q(\omega)=Q(J(\omega),O(\omega))\)
and likewise for \(Z(\omega)\).

The winning superposition \(h_J\) is fixed, while \(h_{J,Q}\)
depends on \(Q\).  On every retained endpoint site, the buffered
property and the transport order fixed after the winning map give
\begin{equation}
 |\mathfrak a_O(h_J)-\mathfrak a_O(h_{J,Q(J,O)})|
       \le C_{\rm tr}\eps^{A_{\rm tr}}\|f\|_2 .
 \label{eq:repayment-winning-candidate-tail}
\end{equation}
The coefficient and count comparison in the original construction
is still available: the squared summed tail has exponent beyond
\(B+1\) after the \(Q_0\) count loss.  For the score we use the
stronger information that the winning amplitude itself has the
floor \eqref{eq:repayment-winner-floor}.  Set
\[
 q_{\rm tr}=
 \frac{C_{\rm tr}S_\eps}{1-q_0}
       \eps^{A_{\rm tr}-A_0/2}.
\]
The choices \(B>2A+3Q_0+10\), \(A>A_0/2+1\), and
\eqref{eq:repayment-transport-order}, together with the subpower
bound on \(S_\eps\), make \(q_{\rm tr}\to0\) at fixed finite
parameters.  Choose \(\eps\) small enough that \(q_{\rm tr}<1/2\).
Equations \eqref{eq:repayment-winner-floor} and
\eqref{eq:repayment-winning-candidate-tail} then imply
\begin{equation}
 m_O(h_J)\le(1-q_{\rm tr})^{-2}m_O(h_{J,Q(J,O)}).
 \label{eq:repayment-winning-candidate}
\end{equation}
In particular \(m_O(h_{J,Q(J,O)})>0\) and \(N_{J,Q(J,O)}>0\) on
every retained sample.  Replacing the winning mass by this
candidate-truncated mass costs at most
\begin{equation}
 3\logeps{1/(1-q_{\rm tr})}
 \label{eq:repayment-candidate-score-cost}
\end{equation}
in the score's mass term.  This cost is included in \(\eta\).

For a fixed \(J,Q\), define
\begin{equation}
 m_Z=\sum_{\substack{O\mapsto Z\\ O\ {\rm in\ the\ localizer}}}
             m_O(h_{J,Q}) .
 \label{eq:repayment-candidate-mass}
\end{equation}
The sum uses every geometric output site mapping to \(Z\) once,
with the allowed bounded lattice multiplicity and its chosen
site-dependent amplitude; it does not sum over copies of a site in
\(\Omega\).  Because \(Z\) includes \(Q\), the entire fiber uses
one input \(h_{J,Q}\).  Thus \(m_Z\) is exactly the
candidate-truncated mass to which the next lemma applies.
The input \(h_{J,Q}\) varies with \(Q\).  These masses are used
only to bound an entropy expression under \(\pi\); they are not
treated as a new admissible packet system for a growth estimate.

\subsection{The repayment estimate}

The prepared cut supplies \(\mu\) and the capacity bound
\eqref{eq:velocity-fiber-regularity}.  The remaining estimate shows
that a near-maximal endpoint score increases the fine-velocity credit.

\begin{proposition}[Finite packet repayment]
\label{prop:finite-repayment}
Fix \(0<\gamma<1\) and a comparison rate
\(b_{\mathrm q}\in[0,1]\).  For the prepared cut constructed above,
let \(\eta\) bound the composition budget \(\mathcal R\),
\(\theta+\xi+(\log C)/\ell\), all logarithmic class, winner, tail,
and label costs, the current mean \(\EE_\pi d_1\), and the finite
classical growth error on the fixed class of lines, after a fixed
increase of \(C_*\).  The selected endpoint and earlier replacement
and composition deficits required in the construction are at most
\(\eta\) on their specified current laws, including the two
conditional prediction marginals.  The common shape and count bins
have width at most \(\eta\); in the isotropic case they also give
the error in \eqref{eq:repayment-isotropic-score}.
When \(k_f>0\), assume \eqref{eq:repayment-selected-credit} under
\(\pi\) with its accumulated error \(\alpha_f\le\eta\).
Finally suppose
\begin{equation}
 P_t\ge B_r+2b_{\mathrm q}h-\Delta .
 \label{eq:repayment-saturation-input}
\end{equation}
Then the prepared statistic \(P,n(P),W,k\) under \(\mu\) satisfies
\begin{equation}
 0\le k\le d h+C_*\eta,\qquad
 k\ge k_f+(4b_{\mathrm q}-2\gamma)h-2\Delta-C_*\eta,
 \quad
 d=\begin{cases}1&\text{elongated},\\2&\text{isotropic}.
                 \end{cases}
 \label{eq:finite-packet-repayment}
\end{equation}
Its pointwise probability bound is
\eqref{eq:repayment-outgoing-credit}.  In the elongated case the
continued aperture normal is within \(C\eps^h\) of \(n(P)\).
The returned score satisfies \eqref{eq:repayment-earlier-scores}
in that case and \eqref{eq:repayment-isotropic-score} in the
isotropic case.  All these bounds are finite at fixed \(\eps\).
\end{proposition}

\subsection{A coherent estimate on the full set of output sites}

We next prove the mass estimate before introducing any positive law
on candidate packets.  The output amplitudes may be chosen separately
at every site.

\begin{lemma}[Packet coherence on a strip or a cell]
\label{lem:packet-coherence}
Use local coordinates with root time \(r\), endpoint duration \(2h\),
and put
\[
 \rho=\eps^u,\qquad s=\eps^{u+h},\qquad
 \sigma=\eps^{u+2h}=s^2/\rho,\qquad
 \lambda_{\rm loc}=s^{-2}.
\]
Fixed multiplicative errors in these relations change only constants.
Let \(P\) be a buffered isotropic cell of width \(\rho\).
Let \(Q\) specify \(P,C_t\) and a position cell of side \(s\) at
\(C_t\).  Candidate atoms have this \(P\), velocity spacing \(s\),
earlier position spacing \(s\), and transported position within
\(Rs\) of that cell, where \(R=\eps^{-\theta}\).
Let \(N_Q\) be their number, and suppose
\[
 |c_T|^2\leq\eps^{-\xi}m.
\]
Partition their velocities into \(K\leq\eps^{-k-\xi}\) bins of
width \(s\), in a fixed normal direction for \(d=1\) or in both
directions for \(d=2\).
Let \(Z\) add to \(Q\) the output position at precision \(\sigma\)
in those \(d\) directions.  Let \(m_Z\) be the sum of squared
packet tests over all output sites mapping to \(Z\), on the input
truncated to these candidates.  Then
\begin{equation}
 \sum_{Z\mid Q}m_Z\leq C\eps^{-\xi}N_Qm,\qquad
 m_Z\leq C\eps^{dh-k-2\theta-2\xi}N_Qm.
 \label{eq:finite-coherence}
\end{equation}
The constant is uniform over all admissible choices of the
site-dependent amplitudes.  Without candidate truncation, the
square root of the mass estimate has the additive tail error
provided by Lemma~\ref{lem:packet-decomposition}.
\end{lemma}

\begin{proof}
The first inequality is the synthesis Bessel estimate followed by
the analysis Bessel estimate of Lemma~\ref{lem:packet-bessel}.
It holds for any subset of output sites and any subset of atoms.

Suppose first that \(d=2\).  If \(A_{OT}\) is the normalized
atom/test matrix, then
\[
 |A_{OT}|\leq C\rho^{-1}\|g_T\|_1
          \leq C s/\rho=C\eps^h.
\]
A full velocity bin contains a bounded number of velocity lattice
centers.  At each such center the proximity condition leaves at
most \(CR^2\) position centers.  Each \(Z\) contains a bounded
number of output sites.  Cauchy--Schwarz first within one bin
and then across its \(K\) bins gives
\[
 m_Z\leq CKR^2\eps^{2h}\sum_T|c_T|^2
       \leq C\eps^{2h-k-2\theta-2\xi}N_Qm.
\]

For \(d=1\), write \(n\) for the bin normal and \(\tau\) for its
orthogonal unit tangent.  Fix one velocity bin \(W\).
There are at most \(CR^2\) candidate atoms per tangent velocity
interval of length \(s\).  We claim
\begin{equation}
 \sum_{O\mapsto Z}\left|
       \sum_{T\in W}c_T A_{OT}\right|^2
 \leq CR^2\eps^h\sum_{T\in W}|c_T|^2.
 \label{eq:one-bin-strip}
\end{equation}

Here are details that also justify using all output sites with
independently chosen amplitudes.  Write \(v=V_P+\rho z\) on a
common bounded patch and surround each output center \(X_O\)
by a square \(B_O\) of side comparable to \(\sigma\).
For \(x\in B_O\), remove the scalar value at \(v=V_P\) from
\[
 \psi_O((v-V_O)/\rho)
 \exp\{i\lambda_{\rm loc}
       [\Phi(C_t,X_O,v)-\Phi(C_t,x,v)]\}.
\]
All derivatives in \(z\) are uniformly bounded, because
\(\lambda_{\rm loc}\rho\sigma=1\).
A common bump and a Fourier series on a fixed enlarged patch
therefore give coefficients \(b_\nu(O,x)\) satisfying
\[
 |b_\nu(O,x)|\leq C_M(1+|\nu|)^{-M}
\]
uniformly in both \(O\) and \(x\).
Weighted Cauchy--Schwarz bounds the squared transform at \(X_O\)
by a fixed rapidly summable combination of squared transforms
at \(x\), each with a common bump and a linear modulation
\(e^{i\nu\cdot(v-V_P)/\rho}\).
These factors are independent of \(x\).
No derivatives with respect to the site index have been used.

Integrate this bound on \(B_O\), divide by its area, and sum.
The squares have bounded overlap, and their union lies in a
strip of normal width \(C\sigma\) and tangent length \(Cs\).
The normalization outside the resulting continuous position
integral is
\begin{equation}
 \rho^{-2}\sigma^{-2}=s^{-4}.
 \label{eq:strip-normalization}
\end{equation}
It is enough to prove the estimate for each of the common
bump/modulation factors, uniformly in its Fourier index; the
factor has bounded modulus and is independent of \(x\).

For fixed normal position, integrate against a smooth majorant
of the tangent interval of length \(Cs\).  For velocities in
the same \(W\), their normal difference is \(O(s)\).  Symmetry
and uniform definiteness of \(\partial_p\zeta\) give
\[
 \partial_\tau\{\Phi(C_t,x,v)-\Phi(C_t,x,v')\}
       =\tau^{\mathsf T}A(x,v,v')(v-v'),
\]
where \(A\) is an average of those definite matrices.
If the tangent separation is a sufficiently large multiple of
\(s\), this derivative has magnitude at least
\(c|(v-v')\cdot\tau|\).
Higher \(x\)-derivatives are bounded by \(C_j|v-v'|\).
Repeated integration by parts in the rescaled tangent variable
thus bounds its kernel by
\[
 C_Ms\left(1+\frac{|(v-v')\cdot\tau|}{s}\right)^{-M}.
\]
For smaller separation the same bound, with another constant,
follows directly by integration over the interval.
Earlier normalized atoms have \(L^1\) norm at most \(Cs\).
Including the normal integration and
\eqref{eq:strip-normalization}, the corresponding Gram entry
is therefore at most
\[
 C_Ms^{-4}\sigma s\,s^2
  \left(1+\frac{|(V_T-V_{T'})\cdot\tau|}{s}\right)^{-M}
 =
 C_M\eps^h
  \left(1+\frac{|(V_T-V_{T'})\cdot\tau|}{s}\right)^{-M}.
\]
The possible overlap of velocity supports changes only the
constant in this estimate.  Sum the rapidly decreasing kernel
over tangent intervals, using the multiplicity \(CR^2\).
The Schur bound proves \eqref{eq:one-bin-strip}.
Finally Cauchy--Schwarz over the \(K\) bins proves the second
inequality of \eqref{eq:finite-coherence} for \(d=1\).
\end{proof}


\subsection{The candidate law and its entropy bound}

We now prove Proposition~\ref{prop:finite-repayment}.
For a value \(J=j\) of positive \(\pi_J\)-probability, let
\(\pi_j\) denote the selected endpoint law conditional on \(J=j\).
Define the finite entropy expression
\[
 \mathcal E_j=
 \Ent_{\pi_j}(Z)+\tfrac12\Ent_{\pi_j}(Z\mid D)
 +\tfrac32\EE_{\pi_j}\logeps{m_Z}
 +\tfrac12\EE_{\pi_j}w_{\rm loc}(D).
\]
All entropies here concern the displayed geometric observations,
rather than the entire sample \(\omega\).  The latter carries
extra labels only so that later augmentation preserves them.
Finite composition first compares the original score with the
conditional winning scores.  Equation
\eqref{eq:repayment-winning-candidate} then replaces their masses
by the candidate-truncated masses.  Finally, conditional log-sum
on each \(O\mapsto Z\) fiber bounds the two entropy--mass terms
by those in \(\mathcal E_j\).  The unconditional fiber calculation
uses all sites mapping to \(Z\); conditional on \(D\), its support
is a subset of those sites.  Thus
\begin{equation}
 P_t\le \Ent_\pi(J)+b_*+\EE_{\pi_J}\mathcal E_J
       +\mathcal R+3\logeps{1/(1-q_{\rm tr})}.
 \label{eq:repayment-candidate-expression}
\end{equation}
This is a comparison of finite expressions under the same endpoint
law.  No growth bound has been applied to the \(Q\)-dependent inputs.

Lemma~\ref{lem:packet-coherence} now gives its two exact bounds for
the masses \(m_Z\) in \eqref{eq:repayment-candidate-mass}.
Only after these analytic inequalities do we introduce a positive
candidate law.  Adjoin an earlier atom to the full endpoint sample by
\begin{equation}
 \widetilde\pi(\omega,T)=\pi(\omega)
   \frac{\boldsymbol1_{\{T\in\mathcal T_{J(\omega),Q(\omega)}\}}}
        {N_{J(\omega),Q(\omega)}} .
 \label{eq:repayment-candidate-law}
\end{equation}
The denominator is positive on every sampled value, as proved above.
This law has exactly the endpoint marginal \(\pi\), including
\(P_f,W_f\) and every stochastic incoming tag.  Conditional on
\((J,Q)\), the atom is uniform on the candidate set and independent
of every other endpoint label.  Its marginal on atoms is generally
biased; it is different from \(\mu\), which is uniform within each
retained class.

For the next calculation fix \(J=j\), and write \(\Ent_j,\EE_j\)
for entropy and expectation under
\(\widetilde\pi(\,\cdot\mid J=j)\).  Put
\[
 B_{1,j}=\EE_j\logeps{N_{J,Q}},\qquad
 S_{t,j}=\Ent_j(C_t\mid T),\qquad
 \mathcal M_j=\EE_j\logeps{m_Z}.
\]
Since \(Z\) includes \(Q\), the sampling rule gives exactly
\[
 \Ent_j(T\mid Z,D)=\Ent_j(T\mid Z)=B_{1,j}.
\]
Candidate proximity gives
\(\Ent_j(Q\mid T,C_t)\le C_*\eta\).  It also gives
\(\Ent_j(Z\mid T,D)\le C_*\eta\).
For the latter bound, \(T,C_t\) locate the position to precision
\(s\), while \(D\) supplies its normal component at precision
\(\sigma\).  In the elongated case the current prediction
\(\angle(n(D),n(P))\le C\eps^{4h}\) moves a residual of size
at most \(\eps^{-\theta}s\) by at most
\(C\eps^{4h-\theta}s< C\sigma\).
In the isotropic case the aperture supplies both fine components.
The remaining costs are the stated lists and boundary-depth errors.

The chain rule gives
\begin{align}
 \Ent_j(Z)&\le \Ent_j(T)-B_{1,j}+S_{t,j}
                 +\Ent_j(Z\mid Q)+C_*\eta,
 \label{eq:repayment-entropy-chain-one}\\
 \Ent_j(Z\mid D)+B_{1,j}
       &\le \Ent_j(T\mid D)+C_*\eta.
 \label{eq:repayment-entropy-chain-two}
\end{align}
Conditional log-sum applied to the aggregate coherence bound, and
the logarithm of its pointwise bound, respectively give
\begin{align*}
 \Ent_j(Z\mid Q)+\mathcal M_j
     &\le B_{1,j}+\logeps{m(j)}+C_*\eta,\\
 \mathcal M_j&\le B_{1,j}+\logeps{m(j)}+k-dh+C_*\eta .
\end{align*}
Use the first inequality with coefficient \(1\) and the second
with coefficient \(1/2\) in
\eqref{eq:repayment-entropy-chain-one}--\eqref{eq:repayment-entropy-chain-two}.
Every occurrence of \(B_{1,j}\) cancels.  The result is
\begin{equation}
 \mathcal E_j\le
 \Ent_j(T)+\tfrac32\logeps{m(j)}
 +\tfrac12\{\EE_jw_{\rm loc}(D)+\Ent_j(T\mid D)
                    +2S_{t,j}+k-dh\}+C_*\eta .
 \label{eq:coherent-entropy-repayment}
\end{equation}
Here the common \(k\) bin absorbs its stated finite bin error.

We assemble these conditional kernels under \(\pi_J\) before
using incoming information.  Averaging
\eqref{eq:coherent-entropy-repayment} gives
\begin{align}
 \EE_{\pi_J}\mathcal E_J
 &\le \Ent_{\widetilde\pi}(T\mid J)
       +\tfrac32\EE_\pi\logeps{m(J)}\notag\\
 &\quad+\tfrac12\bigl\{
       \EE_\pi w_{\rm loc}(D)
       +\Ent_{\widetilde\pi}(T\mid D,J)
       +2\Ent_{\widetilde\pi}(C_t\mid T,J)
       +k-dh\bigr\}+C_*\eta .
 \label{eq:coherent-entropy-repayment-averaged}
\end{align}
The incoming conditional probability bound will be used under this
assembled law, not separately on each \(J\) fiber.

\subsection{A faithful aperture and the incoming information}

The candidate atoms carry exact straight lines through their earlier
positions with their assigned velocities.  In a localizer \(J\),
let \(F_T\) record \(C_t\) and both columns of this line at
isotropic width \(u-h\) at the endpoint.  Its local position
precision is \(\eps^{u+h}=s\).  Thus \(F_T\) is a function of
\((T,C_t,J)\); conditional on \(J\), it need not be a function
of \(T\) alone.

In the elongated case coarsen the orientation of \(D\) and its
aperture data to widths \(0,u-h\), computing the latter from the
center given by \(F_T\).  In the isotropic case use widths
\(u-h,u-h\).  Denote this intermediate observation by \(Y\).
Transport proximity makes \(Y\) determined by \(D\), conditional
on \(J\), up to a list of size at most \(\eps^{-C_*\eta}\).

Enlarge each rectangle of \(Y\) by one fixed factor so that it
contains the entire corresponding isotropic \(F_T\) cell.  The
factor works for every orientation: projections of a square have
length at most its diagonal, and rounding the center adds only a
fixed grid spacing.  This costs at most \((\log C)/\ell\) in depth.
Call the enlarged label \(\bar Y\).  It is a function of \(Y\),
so it has the same conditional list bound from \(D\).
Let \(\kappa(g\mid F_T,J)\) be the conditional law of \(\bar Y\)
under \(\widetilde\pi\), and adjoin \(G\) by that kernel:
\[
 \widetilde\pi(\omega,T,G)
   =\widetilde\pi(\omega,T)\kappa(G\mid F_T,J).
\]
We keep the same notation for the augmented law.  It preserves the
entire old law of \((\omega,T)\) and the joint law of
\((F_T,J,\bar Y)\) with \(\bar Y\) replaced by \(G\).
The new \(G\) is independent of \(T\) conditional on \((F_T,J)\),
contains \(C_t\), and is faithful for every line with that
\(F_T,J\) value.  No \(T\)-dependent cell is appended to its label.

Apply Theorem~\ref{thm:classical-growth} in each localizer with
duration \(2h\) and horizon \(u+h\).  The strict eccentricity
condition, or \(b_*=0\) with \(u_t>h\), gives \(u>h\).
Its root caps have widths at most
\(u+h=u_r-b_*\), and the endpoint horizon \(a+t\le L\), with its
stated depth error, also places them inside the local packet menu.
For every such root test \(F\), the support upper bound gives
\[
 \Ent_j(T\mid F)+\EE_j w(F)\le g(j)+C_*\eta.
\]
This uses only the retained support, so it holds for the biased
candidate marginal.  The positive root supremum is consequently at
most \(g(j)-\Ent_j(T)+C_*\eta\).  Positive growth yields, on
each \(J=j\) fiber,
\begin{equation}
 \EE_j w(G)+\Ent_j(T\mid G)+2S_{t,j}
       \le g(j)+2\gamma h+C_*\eta .
 \label{eq:faithful-classical-bound}
\end{equation}
Averaging these inequalities under the same \(\pi_J\) as before,
\begin{equation}
 \EE_{\widetilde\pi}w(G)+\Ent_{\widetilde\pi}(T\mid G,J)
       +2\Ent_{\widetilde\pi}(C_t\mid T,J)
 \le \EE_\pi g(J)+2\gamma h+C_*\eta .
 \label{eq:faithful-classical-bound-averaged}
\end{equation}

All information in the remainder of this subsection is under the
assembled augmented law \(\widetilde\pi\).  Both \(D\) and \(G\)
contain \(C_t\), and \(F_T\) is a function of \((T,C_t,J)\).
The common time information cancels in the exact chain rule
\begin{align}
 \Mut(T;D\mid J)-\Mut(T;G\mid J)
 &=\Mut(F_T;D\mid C_t,J)-\Mut(F_T;G\mid C_t,J)
       \notag\\
 &\quad+\Mut(T;D\mid F_T,J)-\Mut(T;G\mid F_T,J).
 \label{eq:time-correct-information-splitting}
\end{align}
The last term is zero by the resampling kernel.  Data processing
for \(\bar Y\), followed by preservation of its joint law with
\((F_T,J)\), bounds the first line below by \(-C_*\eta\).
Therefore
\begin{equation}
 \Mut(T;D\mid J)-\Mut(T;G\mid J)
       \ge \Mut(T;D\mid F_T,J)-C_*\eta .
 \label{eq:information-before-credit}
\end{equation}

When \(k_f>0\), the preserved endpoint marginal supplies an
averaged information lower bound
\begin{equation}
 \Mut(T;D\mid F_T,J)\ge k_f-C_*\eta .
 \label{eq:credit-information}
\end{equation}
Here is the finite-list verification.  In absolute units,
\(F_T\) and \(P_f\) have identical velocity depth \(u_t-h\)
and position depth \(t+u_t-h\).  The candidate position differs
from the observed one by at most
\(\eps^{-\theta}\eps^{r+u_r}\), exactly an
\(\eps^{-\theta}\) multiple of that position unit, and the
velocity difference is \(O(\eps^{u_t})\).
The localizer has the coarser width in
\eqref{eq:repayment-localizer-depth}; its remaining tags have the
charged conditional count.  It follows that
\[
 \Ent(F_T,J\mid P_f)\le C_*\eta,\qquad
 \Ent(P_f\mid F_T,J)\le C_*\eta .
\]
Because \(\widetilde\pi_\Omega=\pi\),
\eqref{eq:repayment-selected-credit} still holds for its endpoint
labels.  It is not asserted after conditioning separately on \(J\).
The chain rule gives
\[
 \Ent(W_f\mid F_T,J)
 \ge \Ent(W_f\mid P_f)-\Ent(F_T,J\mid P_f)
 \ge k_f-C_*\eta .
\]

Neither coarse label alone determines \(W_f\).  The pairs
\((D,F_T,J)\) and \((T,F_T,J)\), however, each determine a
list of at most \(\eps^{-C_*\eta}\) possibilities.
First list the possible \(P_f\)'s supplied by \(F_T,J\).
For each listed value, \(n_f(P_f)\) is known, with no continuity
assumption on this function.  Use that value's coarse velocity center
as reference.  The remaining velocity has size
\(O(\eps^{u_t-h})\).  The incoming alignment of the normal of
\(D\) with \(n_f(P_f)\) is \(O(\eps^h)\), so changing between
them changes the measured component by \(O(\eps^{u_t})\), one
\(W_f\) unit.  Thus \(D\) resolves that component to the stated
list.  The exact velocity of \(T\) does so as well, since it differs
from the output velocity by \(O(\eps^{u_t})\).
For isotropic tests use both velocity coordinates.  Consequently
\[
 \Ent(W_f\mid D,F_T,J)+\Ent(W_f\mid T,F_T,J)\le C_*\eta .
\]
The elementary information inequality
\[
 \Mut(T;D\mid F_T,J)\ge
 \Ent(W_f\mid F_T,J)-\Ent(W_f\mid D,F_T,J)
                         -\Ent(W_f\mid T,F_T,J)
\]
proves \eqref{eq:credit-information}.  This is a conditional
information averaged over \(J\) under \(\widetilde\pi\), which is
exactly what the averaged bounds need.  If \(k_f=0\), use
nonnegativity of conditional mutual information instead.

Combining \eqref{eq:information-before-credit} and
\eqref{eq:credit-information} with
\eqref{eq:faithful-classical-bound-averaged} gives
\[
 \EE_{\widetilde\pi}w(G)+\Ent_{\widetilde\pi}(T\mid D,J)
       +2\Ent_{\widetilde\pi}(C_t\mid T,J)
 \le \EE_\pi g(J)+2\gamma h-k_f+C_*\eta .
\]
The weight change on each localizer is
\[
 w_{\rm loc}(D)-w(G)=
 \begin{cases}
 (1-\gamma)h+O(\eta),&d=1,\\
 2h+O(\eta),&d=2,
 \end{cases}
 \qquad
 w_{\rm loc}(D)-w(G)\le dh+C_*\eta .
\]
Also \(\Ent_{\widetilde\pi}(T\mid J)\le
\EE_\pi\logeps{N(J)}\), since every candidate lies in its retained
class.  Insert these bounds into
\eqref{eq:coherent-entropy-repayment-averaged}, then use
\eqref{eq:repayment-candidate-expression}.  The tail comparison
costs are included in \(\eta\), so after returning to the original
units,
\[
 P_t\le B_r+\gamma h+\tfrac12(k-k_f)+C_*\eta .
\]
Comparison with \eqref{eq:repayment-saturation-input} proves the
lower bound in \eqref{eq:finite-packet-repayment}.  Its upper
bound is \eqref{eq:velocity-fiber-regularity}.

Under \(\mu\) at time \(r\), \(P\) has depth \(u_t=u_r-h\)
and \(W\) has depth \(u_r\).  Thus
\eqref{eq:repayment-outgoing-credit} and the earlier aperture and
score bounds supply the next cut's incoming data.  That cut may use
a new atom decomposition: the coarse-label lists carry the observable
\(W\mid P\) bound, without identifying the two temporary candidate
populations.
This completes Proposition~\ref{prop:finite-repayment}.


\subsection{A finite backward chain}

Suppose, for a contradiction, that
\(\betaq>\gamma/2\), and set
\[
 c=4\betaq-2\gamma>0.
\]
Choose an integer \(q\) with \(qc>2\).
Proposition~\ref{prop:boundary-alternatives} supplies \(q\)
backward steps of a common positive length \(2h\), with
all times positive and \(u_t>h\).
In the forced alternative first take \(h\) sufficiently small
compared with the fixed aspect gap times \(\gamma\);
this verifies \eqref{eq:repayment-eccentricity}.
In the other alternative all the back tests are exactly isotropic.

Fix this chain and start with \(k_0=0\).  At cut \(i\), apply the
rule for a prepared cut and call its final endpoint law \(\pi_i\);
the returned \(\mu_i\) is the next endpoint law.  Let
\(\Delta_i,\eta_i\) be the actual selected errors for \(\pi_i\),
with \(J\)-averages under \((\pi_i)_J\).  They include the exceptional-fiber score splits and
the accumulated \(\zeta+\logeps{1/p_A}\) with every \(p_A\)
measured in its current law.  Choose the finitely many retention
parameters and current-law branches so this finite cost sum is
as small as required.

Within each fixed outer row \(j\), choose the finite inner parameters
\(\theta,\xi\) and the finitely many \(\zeta\)'s, then the initial
packet tail orders and coefficient cutoff, the later transport order,
and finally the precision \(\eps\).  These choices make the inner
errors, including both mass-comparison costs, arbitrarily small in
that row.  Finite composition and
\eqref{eq:repayment-input-bridge} leave the selected deficits bounded
by that row's \(\Delta_j\), with the fixed finite selection factors,
plus those inner errors.  With the target shape and \(g\) accuracies
fixed before choosing the late row, these factors are bounded across
the subsequence: only uniformly positive current masses of qualifying
unions and retained complements divide deficits, while individual
tag probabilities enter only logarithmic costs.
The selected deficits are not sent to zero inside the row.
Only then choose a sufficiently late outer row, with \(q\) and \(h\)
already fixed across the subsequence, and choose finite inner precision
in that row so that the total error bound below holds.  No transport
order uniform in unbounded outer aspect parameters is required.

Apply Proposition~\ref{prop:finite-repayment} with
\(b_{\mathrm q}=\betaq\).
Iteration of \eqref{eq:finite-packet-repayment} gives
\[
 k_q\geq qch-\sum_{i=1}^q(2\Delta_i+C_*\eta_i),
 \qquad
 k_q\leq2h+C_*\eta_q.
\]
Choose the finite errors so that
\[
 \sum_{i=1}^q(2\Delta_i+C_*\eta_i)+C_*\eta_q
                  <(qc-2)h.
\]
The two inequalities are incompatible.
Consequently \(\betaq\leq\gamma/2\), completing the proof of
Theorem~\ref{thm:packet-growth}.
